\begin{table}[H]
  \centering
  \caption{The four statements of Part~II.  All are conditional schemas
  with primitive predicates; none is proved.}
  \label{tab:recognition:summary}
  \small
  \renewcommand{\arraystretch}{1.1}
  \begin{tabularx}{\textwidth}{@{}>{\raggedright\arraybackslash}p{0.27\textwidth}>{\raggedright\arraybackslash}X@{}}
    \toprule
    Statement & Content \\
    \midrule
    \Cref{thm:template-applicability} & On a trace-state-only formed closure satisfying the SB closure assumption, a seven-stage record for which the template applies yields theorem-grade landing. \\
    \Cref{thm:verdict-signature} & Under the same hypotheses, the derivation sweep returns the verdicts \SigPattern. \\
    \Cref{thm:closure-content-thesis} & The load-bearing content of a trace-state-only closure is formed-layer content, is not derivable from the primitives, and is named as an accepted source. \\
    \Cref{thm:attack-foreclosure-v5} & For an in-scope target and a trace-state-only formed closure satisfying the SB closure assumption: derivation routes do not advance the target without outside content, and recognition-mode landings import the load-bearing content and land at theorem grade. \\
    \bottomrule
  \end{tabularx}
\end{table}
